\documentclass[11pt]{article}

\usepackage[margin=1in]{geometry}
\usepackage[T1]{fontenc}
\usepackage[utf8]{inputenc}
\usepackage{lmodern}
\usepackage{setspace}
\usepackage{hyperref}
\usepackage{csquotes}
\usepackage[authoryear]{natbib}

\title{Berkeley as a Precursor to Embodied Machine Intelligence:\\
Symbolic Logic, Control, and Behavior Before Adaptive Robotics}
\author{}
\date{\today}

\begin{document}

\maketitle

\begin{abstract}
Edmund C. Berkeley is usually remembered as a writer who helped connect symbolic logic to computing machinery. That description is correct, but incomplete. This paper argues that Berkeley's \emph{Symbolic Logic and Intelligent Machines} should also be read as an early contribution to embodied machine intelligence. Berkeley does not merely treat intelligence as formal symbol manipulation. He repeatedly defines intelligent machines in operational terms, ties logical form to hardware realization, and describes intelligent behavior as the coordinated interaction of inputs, outputs, memory, calculation, and control. Most importantly, his treatment of robots, machine activities, and the algebra of states and events moves beyond static logical form toward temporally extended, environment-coupled behavior. Berkeley does not yet reach the later adaptive and reinforcement-oriented framework found in David L. Heiserman's evolutionary machine intelligence, but he does mark an important intermediate position between algebraic logic and creature-like machine behavior. Read in this way, Berkeley appears not only as a symbolic logician of machines, but as a bridge figure in the history of embodied machine intelligence.
\end{abstract}

\onehalfspacing

\section{Introduction}

Edmund C. Berkeley is commonly placed in the history of symbolic logic, early computing, and machine thought. That placement is justified. \emph{Symbolic Logic and Intelligent Machines} is overtly concerned with Boolean algebra, logical relations, class reasoning, and the design of information-handling machinery \citep{berkeley1959symbolic}. Yet this standard placement risks narrowing Berkeley's significance. He was not only interested in the formal calculability of statements and classes. He was also interested in how reason-like competence could be built into machinery that senses, stores, decides, and acts.

The aim of this paper is to argue that Berkeley should be read as a bridge figure between symbolic logic and embodied machine behavior. He is not a full theorist of adaptive embodied AI in the later Heiserman sense, but neither is he adequately understood as a purely formal or disembodied logician. His work gives a logically rigorous account of how physically organized machines can take in information, retain it, operate on it, and guide their own activity over time \citep[pp.~66--70]{berkeley1959symbolic}.

\section{Research Question and Thesis}

The guiding question is how Berkeley's \emph{Symbolic Logic and Intelligent Machines} should be situated in the history of machine intelligence: as a work of symbolic logic applied to machinery, or as an early account of embodied, behavior-organizing intelligence.

The working thesis is that Berkeley's work is best read as both, but with stronger embodied implications than standard histories usually acknowledge. He inherits formal machinery from Boole, Venn, and Shannon, yet uses that inheritance to theorize intelligent systems as physically realized, temporally organized, and behaviorally active machines \citep{boole1847mathematical,boole1854laws,venn1894symbolic,shannon1938symbolic}. His account stops short of adaptive creature intelligence, but it clearly anticipates later embodied and control-oriented approaches \citep{heiserman1981robot}.

\section{Literature Review}

The literature relevant to this paper falls into five overlapping bands. The first and most important consists of the primary sources that directly carry the argument: Berkeley, Boole, Venn, Shannon, and Heiserman. Berkeley is the principal object of interpretation. Boole provides the originating algebraic logic that Berkeley explicitly treats as foundational \citep{boole1847mathematical,boole1854laws}. Venn represents the strengthened class-logic and symbolic tradition that Berkeley inherits and occasionally names directly \citep{venn1894symbolic}. Shannon supplies the crucial bridge from logical relations to relay and switching hardware \citep{shannon1938symbolic}. Heiserman, by contrast, is not part of Berkeley's own lineage, but functions as a later comparison point that clarifies what Berkeley does and does not achieve in relation to embodied and adaptive machine intelligence \citep{heiserman1981robot}.

The second band consists of immediate contextual works in symbolic logic, especially Basson and O'Connor's \emph{Introduction to Symbolic Logic}. These are useful not because they define Berkeley's project, but because they help establish what a more orthodox mid-century symbolic-logic presentation looked like \citep{basson1960introduction}. Against that background, Berkeley appears less like a standard textbook expositor and more like a writer pressing logic toward hardware, behavior, and machine organization.

The third band is the cybernetic and control literature: Wiener, Ashby, Turing, and Walter \citep{wiener1948cybernetics,ashby1952design,ashby1956introduction,turing1948intelligent,turing1950computing,walter1953living}. These works are not needed to prove Berkeley's core claims, but they help situate him within a broader postwar field concerned with control, communication, regulation, and autonomous machinery. Their function in this paper is therefore contextual and comparative, not genealogical.

The fourth band is later behavior-based or embodied-machine comparison, especially Braitenberg and Brooks \citep{braitenberg1984vehicles,brooks1986robust,brooks1991intelligence}. These references are used cautiously. They do not show influence on Berkeley. Instead, they show that later discussions of intelligent behavior often shifted away from symbolic representation toward sensorimotor organization and situated control. Berkeley becomes interesting in this frame because he appears as an earlier figure who still retains logical structure while nonetheless treating machines as active, organized, and environment-facing systems.

The fifth and lightest band is secondary historical or philosophical interpretation. This paper uses such literature only sparingly. Its method is primarily analytic and primary-source driven. The goal is not to subsume Berkeley under an already settled historiography, but to show from the texts themselves that his work occupies a more behavior-centered and embodied position than is usually granted.

\section{Logic Embodied in Machinery}

One of Berkeley's major contributions is to insist that symbolic logic is not merely a way of describing reasoning in abstraction. It is also a way of designing machinery. This is the decisive move from logic as formal analysis to logic as implemented intelligence. The historical bridge is Shannon. Once Boolean relations can be realized in relay and switching circuits, conditions and inferences no longer remain on paper or in thought alone. They become structures of hardware \citep{shannon1938symbolic}.

Berkeley adopts this possibility as fundamental. In his treatment, classes, statements, and logical conditions become design elements for circuits, control systems, and automatic machines. This means that logic is not simply explanatory; it is operational. Berkeley's interest in symbolic logic is therefore inseparable from his interest in machinery capable of carrying out organized, reason-like operations in the world \citep{berkeley1959symbolic}.

\section{Intelligent Machines as Organized Systems}

The strongest evidence for an embodied reading appears in Berkeley's own definitions and examples. He lists a wide range of machines as ``intelligent machines,'' including traffic-light controllers, air-traffic systems, factory control equipment, game-playing machines, reading and recognizing devices, robots, and digital computers \citep[pp.~66--69]{berkeley1959symbolic}. These are not merely disembodied theorem provers. Many are directly engaged with environments, signals, movement, and control.

Berkeley then proposes a general architecture shared by almost every intelligent machine. Such a machine, he says, has inputs, outputs, storage or memory, calculating capacity, and a control unit \citep[p.~70]{berkeley1959symbolic}. This matters conceptually because it frames intelligence not as a single faculty but as the coordination of sensing, acting, remembering, processing, and directing. Especially important is Berkeley's remark that the control unit may itself be subject to instructions calculated by the machine, allowing the machine in some sense to guide itself \citep[p.~70]{berkeley1959symbolic}. This is modest autonomy, but it is autonomy nonetheless.

\section{States, Events, and Temporally Extended Behavior}

If Berkeley had stopped with Boolean algebra and machine architecture, the case for an embodied reading would remain suggestive. It becomes much stronger in his treatment of the algebra of states and events. There Berkeley explicitly argues that Boolean calculus must be extended in order to come ``into closer relation with the real world,'' since classes and statements change with time and since machine operation involves changing states and events rather than static truths alone \citep[pp.~145--147]{berkeley1959symbolic}.

A state lasts for a time; an event is brief. This distinction allows Berkeley to formalize persistence, interruption, transition, delay, and sequence. He also insists on the practical importance of time units, step-functions, and even intervals in which a machine state is temporarily undefined during transition \citep[pp.~145--147]{berkeley1959symbolic}. The importance of this move is difficult to overstate. Intelligent behavior becomes describable as a structured passage through temporally extended conditions rather than as a mere arrangement of timeless logical forms. In this respect Berkeley moves from logic as static relation toward logic as an account of organized activity.

\section{Robots and Behavioral Programs}

Berkeley's robot chapters make the behavioral implications even clearer. He defines a robot as a machine able to ``behave'' by itself, taking in sensations by means of sensing devices, performing actions by means of acting devices, and correlating sensations and actions by means of circuits or mechanisms expressing instructions \citep[pp.~177--178]{berkeley1959symbolic}. This is already much closer to an agent-based conception of intelligence than to a purely symbolic one.

His examples are revealing. Berkeley's electronic robot squirrel ``Squee'' is organized around activities such as hunting, homing, and depositing, each of which persists until relevant conditions terminate it \citep[pp.~177--180]{berkeley1959symbolic}. Likewise, the ``robot animal'' problem is introduced explicitly as a programming problem in which one must assign a number of states to a robot and then specify changes from one state to another when certain events happen \citep[pp.~177, 180--181]{berkeley1959symbolic}. These are not merely symbolic demonstrations. They are behavior-organizing schemes in which logical conditions select and govern activity.

\section{Berkeley and Heiserman}

The comparison with David Heiserman clarifies both Berkeley's importance and his limits. Heiserman's \emph{Robot Intelligence ... with Experiments} presents an evolutionary hierarchy in which Alpha-class creatures rely on constrained reflexive activity, Beta-class creatures add a memory of past successful responses, and Gamma-class creatures generalize from prior experience while building confidence in successful patterns of action \citep[pp.~18--20]{heiserman1981robot}. In that framework, machine intelligence is lived by a creature in an environment. It senses, acts, remembers, adapts, and in the Gamma case anticipates.

Berkeley does not reach this level of adaptivity. His machines are mostly engineered top-down rather than self-modifying through reinforcement or confidence-based generalization. He does not provide Beta-style stored response schemes or Gamma-style conjectural extensions of prior success. Even so, the continuity remains important. Berkeley already supplies the control-and-behavior side of the story: intelligence as physically realized organization, behavior as transitions among activities, sensing and acting as core capacities, and internal control as the basis of outward competence.

The difference, then, is not that Berkeley lacks embodiment. The difference is that Berkeley's embodiment is primarily engineered and prescribed, whereas Heiserman's becomes increasingly adaptive and experiential. Berkeley occupies the earlier, intermediate position: he links logic and circuitry to organized machine behavior before later writers extend that connection into creature-like adaptation.

\section{Conclusion}

Berkeley should not be reduced to a symbolic logician who happened to mention machines. His work deserves to be read as an effort to show how intelligence can be built into physically organized systems that sense, remember, control, and act over time. His taxonomy of intelligent machines, his five-part architectural scheme, his extension of logic into states and events, and his robot examples together support a reading of Berkeley as an early contributor to embodied machine intelligence in a control-oriented sense.

This conclusion should be stated carefully. Berkeley does not formulate adaptive, reinforcement-based, or homeostatic machine intelligence in the later Heiserman manner. But that limitation clarifies his historical importance rather than diminishing it. Berkeley occupies a crucial intermediate position between algebraic logic and adaptive creature intelligence. He shows how formal logical structure can govern temporally extended, physically embodied activity. In that sense he is not the endpoint of embodied machine intelligence, but he is clearly part of its prehistory.

\bibliographystyle{plainnat}
\bibliography{references}

\end{document}
